\documentclass{article}

% On Overleaf, switch to the NeurIPS 2026 template's style (the course's template) and keep the same option.
\usepackage[preprint]{neurips_2025}

\usepackage[utf8]{inputenc}
\usepackage[T1]{fontenc}
\usepackage{hyperref}
\usepackage{url}
\usepackage{booktabs}
\usepackage{amsfonts}
\usepackage{amsmath}
\usepackage{nicefrac}
\usepackage{microtype}
\usepackage{xcolor}
\usepackage{graphicx}

% Placeholders for numbers that come from the evaluation run; none may remain in the submitted version.
\newcommand{\todo}[1]{\textcolor{red}{[TODO: #1]}}

\title{LA-Airbnb: A Text-Tabular Price Prediction Dataset\\Curated with the MulTaBench Pipeline}

\author{
  Diyar Hasian \\
  Technion -- Israel Institute of Technology \\
  \texttt{hasian.diyar@gmail.com}
}

\begin{document}

\maketitle

\begin{abstract}
Tabular datasets whose features include free text are common in practice, yet few public benchmarks verify that
their text actually matters for prediction. We curate \emph{LA-Airbnb}, a regression dataset of 22,569 Los Angeles
Airbnb listings that pairs 47 structured features with six free-text fields written by hosts, and predict the
logarithm of the nightly price. We evaluate it with the MulTaBench curation pipeline: five tabular learners (TabM,
CatBoost, LightGBM, TabPFN~2.0 and TabPFN~2.5), each under four conditions (structured features only, text only,
both with frozen E5 embeddings, and both with target-aware LoRA-tuned embeddings). \todo{one sentence with the verdict:
how many learners satisfy joint signal and TAR gain, and the average gains}. Data, code and per-run results are public.
\end{abstract}

\section{Introduction}

Many real-world tables mix modalities: next to numerical and categorical columns, they hold free text such as product
descriptions, clinical notes or reviews. Standard tabular learners---gradient-boosted decision trees
\cite{chen2016xgboost,prokhorenkova2018catboost,ke2017lightgbm} and tabular foundation models
\cite{hollmann2025accurate,grinsztajn2025tabpfn25,qu2025tabicl}---cannot read raw text, so text is typically encoded
by a pretrained sentence encoder such as E5 \cite{wang2022e5} and appended as numerical features. Joint models such as
TabSTAR \cite{arazi2025tabstar} go further and adapt the text encoder to the prediction target.

Comparing such methods requires datasets in which both modalities carry complementary signal, and in which the task
rewards adapting the text representation. MulTaBench \cite{arazi2026multabench} formalizes these requirements as a
\emph{curation pipeline}: a dataset is admitted only if joint models beat both unimodal baselines (\emph{joint signal})
and target-aware representations beat frozen ones (\emph{TAR gain}). Many candidate datasets fail: in particular, the
Airbnb datasets already in MulTaBench (Melbourne, Seattle and New York) appear only in its extended tier, which
requires joint signal alone.

We contribute \emph{LA-Airbnb}, a cleaned text-tabular dataset built from the Inside Airbnb scrape of Los Angeles
\cite{insideairbnb}, together with a documented cleaning procedure and a full evaluation under the MulTaBench
protocol (5 learners $\times$ 4 conditions $\times$ 6 folds). \todo{what the project achieves: passes or not, and
the headline deltas}.

\section{Related Work}

\textbf{Tabular learners.} Gradient-boosted trees, including XGBoost \cite{chen2016xgboost}, CatBoost
\cite{prokhorenkova2018catboost} and LightGBM \cite{ke2017lightgbm}, remain strong defaults on tabular data. TabM
\cite{gorishniy2024tabm} is an efficiently ensembled MLP. Tabular foundation models such as TabPFN
\cite{hollmann2023tabpfn,hollmann2025accurate}, TabPFN~2.5 \cite{grinsztajn2025tabpfn25} and TabICL
\cite{qu2025tabicl} predict in-context after pretraining on synthetic tables. All of these are unimodal: they consume
text only through precomputed embeddings.

\textbf{Text-tabular learning.} The common recipe embeds each text column with a frozen encoder like E5
\cite{wang2022e5}, compresses it and concatenates it with the structured features. TabSTAR
\cite{arazi2025tabstar} instead learns target-aware text representations; parameter-efficient fine-tuning with LoRA
\cite{hu2022lora} makes such adaptation cheap.

\textbf{Benchmarks.} MulTaBench \cite{arazi2026multabench} collects 80 text- and image-tabular datasets, of which 40
form its core after passing both curation criteria. Our work applies its pipeline to a new dataset.

\section{The LA-Airbnb Dataset}

\textbf{Source.} We start from the Inside Airbnb scrape of Los Angeles of 2--4 May 2017 \cite{insideairbnb}, as
distributed on Kaggle: 31,253 listings with 95 columns. Los Angeles is not among MulTaBench's datasets.

\textbf{Columns.} We keep six free-text fields written by the host: the listing title (\texttt{name}, median 6
words), \texttt{description} (164), \texttt{neighborhood\_overview} (47), \texttt{transit} (32),
\texttt{house\_rules} (43) and \texttt{host\_about} (49). The \texttt{summary} and \texttt{space} fields are dropped,
since \texttt{description} contains them in 97\% and 99.7\% of listings respectively. The 47 structured features
cover the property (type, room type, capacity, bedrooms, bathrooms, beds, amenity count), location (neighbourhood,
latitude, longitude), host (tenure, response time and rate, Superhost status, verifications, number of listings),
booking terms (fees, minimum and maximum nights, cancellation policy, instant booking), availability and reviews
(counts, recency and the seven review scores). Dates become days before the scrape and prices become numbers;
categories with fewer than 20 listings are merged. We drop identifiers, URLs, host names, constant or empty columns,
redundant location fields (street, city, zip code), and mask e-mail addresses and phone numbers in the text.

\textbf{Row filters.} We drop listings that were never booked (7,200; their asking price is unvalidated, and they
include all \$10,000-per-night listings), whose description has fewer than 15 words (1,163) or is mostly in a
non-Latin script (741; E5-small-v2 is an English encoder), priced under \$20 (144) or above \$1,000 per guest
(24, e.g.\ \$9,000 for a one-bedroom apartment), requiring stays over 30 nights (169), or listing no guests or
bedrooms (35); filters overlap, removing 8,177 listings in total. Finally, we drop 492 listings whose description
duplicates another listing's (mostly one host's near-identical units), which would otherwise appear in both train and
test splits. The result has 22,569 listings.

\textbf{Target.} We predict $\log(\text{price})$ per night. Prices span \$20--\$10,000 (median \$100) with skewness
12.9, and the top 1\% of listings hold 72\% of the variance; on the log scale the skewness is 0.8 and errors are
relative, so every listing contributes comparably to the RMSE. Weekly and monthly prices, set from the nightly price,
are removed.

\section{Experimental Setup}

We use the official MulTaBench code (commit \texttt{ff0c9c8}) and its protocol: three-fold cross-validation repeated
twice with a fixed seed (6 folds), each fold capped at 10,000 training and 5,000 test rows, and default
hyperparameters for TabM, CatBoost, LightGBM, TabPFN~2.0 and TabPFN~2.5. Each learner runs under four conditions:
\emph{Structured} (the 47 structured features), \emph{Unstructured} (the six text fields only), \emph{Joint Frozen}
(all features, text embedded by frozen E5-small-v2) and \emph{Joint TAR} (all features, text embedded by E5-small-v2
fine-tuned on the target). Each text column is embedded as ``\texttt{passage: <column>: <text>}'', standardized
and reduced to 30 principal components fitted on the training rows. For TAR, LoRA \cite{hu2022lora} (rank 16) tunes
the last three layers of one shared E5 on every text column of the fold's training rows, with the target
discretized into 20 bins, for up to 50 epochs with early stopping on a validation split; as in MulTaBench, one tuned
encoder per fold serves all five learners. The score is MulTaBench's negative RMSE on $\log(\text{price})$ (higher
is better); we also report $R^2$. Following the course's criteria, a learner passes if Joint Frozen strictly beats
both unimodal conditions and Joint TAR strictly beats Joint Frozen, on the mean over folds; the dataset passes if at
least three of the five learners do.

\section{Results}

\todo{Table~\ref{tab:results}: mean $R^2$ (and RMSE) $\pm$ std over 6 folds, per learner and condition, from
results/la\_airbnb\_price/evaluation\_report.csv}

\begin{table}[h]
  \caption{$R^2$ on $\log(\text{price})$, mean $\pm$ std over 6 folds. Bold: best condition per learner.}
  \label{tab:results}
  \centering
  \begin{tabular}{lcccc}
    \toprule
    Learner & Structured & Unstructured & Joint Frozen & Joint TAR \\
    \midrule
    TabM        & \todo{} & \todo{} & \todo{} & \todo{} \\
    CatBoost    & \todo{} & \todo{} & \todo{} & \todo{} \\
    LightGBM    & \todo{} & \todo{} & \todo{} & \todo{} \\
    TabPFN 2.0  & \todo{} & \todo{} & \todo{} & \todo{} \\
    TabPFN 2.5  & \todo{} & \todo{} & \todo{} & \todo{} \\
    \bottomrule
  \end{tabular}
\end{table}

\todo{Curation verdict per learner (from curation\_verdict.csv): joint signal, TAR gain, their margins and on how
many of the 6 folds each held. State whether the dataset passes (at least 3 of 5).}

\todo{Analysis: how much text adds over structured features (Joint Frozen $-$ Structured), how much the structured
features add over text, and how large the TAR gain is relative to fold-to-fold variation.}

\section{Conclusion}

We curated LA-Airbnb, a text-tabular price prediction dataset of 22,569 listings with six host-written text fields,
and evaluated it with the MulTaBench curation pipeline across five learners and four conditions.
\todo{summary of the outcome}.

\textbf{Limitations.} The data covers one city at one point in time (May 2017). Like MulTaBench, we split IID, so
hosts with several listings can appear in both train and test splits, which may overstate performance, although
removing duplicated descriptions mitigates the most direct overlap. Asking prices are not transaction prices, and
the filters on reviews and language shift the population toward established, English-language listings.

\textbf{Future work.} Group-aware splits by host, other targets from the same table (review rating, Superhost
status), other cities and scrapes of Inside Airbnb, and adding the listing photos as a third modality.

\bibliographystyle{plain}
\bibliography{references}

\end{document}
